\section{Conclusions and Future Work}
\label{sec:conclusions}

\subsection{Conclusions}

A discrete Galerkin--POD reduced-order model was developed
for the parameterized steady-state energy equation. The
full-order model was used to generate temperature
snapshots, from which an original POD basis was
constructed. The discretized governing equation was
then projected onto the retained modes to obtain a
reduced system for calculating the modal coefficients.

The original basis was assessed using 100 independent
LHS samples and different retained mode counts.
Increasing the number of modes substantially reduced
the prediction errors. The mean relative $L_2$ error
decreased from 39.56\% with one mode to 0.0468\%
with ten modes.

An adaptive snapshot-enrichment procedure was
subsequently developed. Three high-error parameter
regions were selected, and 15 additional full-order
snapshots were generated around them. These snapshots
were combined with the original 100 snapshots to
construct an augmented POD basis.

The final comparison used 100 test cases excluded
from both POD training datasets. The results showed
that enrichment did not improve the ROM uniformly.
At two to five modes, the mean relative error
increased. At six, eight, and ten modes, it decreased
by 40.77\%, 83.91\%, and 96.98\%, respectively.
At ten modes, the maximum relative error also
decreased by 96.92\%. However, only 20\% of the
individual test cases improved at this mode count.

These findings indicate that targeted snapshot
enrichment can substantially reduce large prediction
errors at higher basis dimensions, but may degrade
accuracy for other parameter combinations. Its
effect therefore depends on both the retained
basis dimension and the location within the
parameter domain.

\subsection{Future Work}

Future work can investigate alternative hotspot
selection criteria and parameter-space sampling
strategies to improve the distribution of prediction
errors. Different enrichment rules and mode-dependent
basis construction could also be examined to obtain
more consistent improvements across test cases.

The evaluation could be extended to larger independent
test datasets to further assess generalization.
Finally, precomputing and storing reusable reduced
operators would improve the offline-online structure
of the implementation and provide a more complete
assessment of its computational efficiency.
